How is the general linear model defined with heteroscedacity?


The model \( y = X \beta + \varepsilon \) 
is called the general linear model if the following assumptions hold:

\( \mathbb{E}(\varepsilon) = 0 \)
\( \operatorname{Cov}(\varepsilon) = \sigma^2 \mathbf{W}^{-1} \)
for some known positive definite matrix \( \mathbf{W} \) 
\( \operatorname{rk} \left( \mathbf{X} \right) = k + 1 = p \leq n \)
If additionally errors are Gaussian, the model is called the general normal regression.
For stochastic covariates all properties are understood conditionally on the design matrix \(\mathbf{X} \).


Explain the significance of the GLS? How are estimators and their properties ad-
justed in comparison to the classical case?

Keep in mind the assumptions of the GLS:
\( y = X \beta + \varepsilon \)
\( \mathbb{E}(\varepsilon) = 0 \)
\( \operatorname{Cov}(\varepsilon) = \sigma^2 \mathbf{W}^{-1} \).


1. Aitken estimator
The weighted least squares estimator can easily be computed by diagonalization:
\( \hat{\beta} = (X^{\top}WX)^{-1}X^{\top}Wy \) 
2. Properties
\( \mathbb{E}(\hat{\beta}) = \beta \)
\( \operatorname{Cov}(\hat{\beta}) = \sigma^2 (X^{\top}WX)^{-1} \)

Gauß-Markov Theorem:
Among all unbiased and linear estimators of \( \beta \) the Aitken estimator has minimal variance.
3. REML-Estimator for \( \sigma^2 \)
The REML-estimator is unbiased:
\( \hat{\sigma}^2 = \frac{1}{n - p} \hat{\varepsilon}^{\top} W \hat{\varepsilon} \)
Remark
All properties do, of course, only hold under the condition that the model is correctly specified.
In particular \(W\) is known, while \( \beta \) and \( \sigma^2 \) are non-random, but unknown model parameters.



Illustrate the GLS for data from groups of entities.

So far, estimation relied on observations \((y_i, x_i)\) with distinct covariate vectors \(x_i, i = 1, 2, \dots, n\).
Suppose now only \( G \) different covariate vectors occur. Then the data can be separated into \( G \) groups with \( n_g \) observations 
and average response \( \bar{y}_g \), \( g = 1, 2, \dots, G \), respectively:

\[
\begin{array}{cccc}
\text{Group 1} & \begin{bmatrix} n_1 \\ \end{bmatrix} & \begin{pmatrix} \bar{y}_1 \\ \end{pmatrix} & \begin{pmatrix} 1 & x_{1,1} & \dots & x_{1,k} \\ \end{pmatrix} \\
\text{Group 2} & \begin{bmatrix} n_2 \\ \end{bmatrix} & \begin{pmatrix} \bar{y}_2 \\ \end{pmatrix} & \begin{pmatrix} 1 & x_{2,1} & \dots & x_{2,k} \\ \end{pmatrix} \\
\vdots & \vdots & \vdots & \vdots \\
\text{Group G} & \begin{bmatrix} n_G \\ \end{bmatrix} & \begin{pmatrix} \bar{y}_G \\ \end{pmatrix} & \begin{pmatrix} 1 & x_{G,1} & \dots & x_{G,k} \\ \end{pmatrix}
\end{array}
\]

This situation can be handled by taking 

\[W = \text{diag} \left( \frac{1}{n_1}, \frac{1}{n_2}, \dots, \frac{1}{n_G} \right)\]



Explain briefly regularization techniques like Ridge regression and LASSO.

Least squares estimation might be unstable.
Possible reasons are:

Design matrix includes columns that are almost collinear.
Number of regression effects \(p\) is large, but \(n\) is comparably small.

Regularization techniques stabilize numerical estimation procedures and may be employed for variable selection.
These are also applicable in the general linear model, but for simplicity we focus on the classical linear model.
We will consider the following techniques:

- Penalized Least Squares Criteria
- Ridge Regression
- LASSO


Ridge regression refers to a regularization that uses the squared \(l^2\)-norm of the estimator as penalty:

\(\text{pen}(\beta) = \beta^\top \beta\)

The minimization problem 
\(\text{PLS}(\beta) = (y - X \beta)^\top (y - X \beta) + \lambda \cdot \beta^\top \beta\)
has the solution
\(\hat{\beta}_{\text{PLS}} = (X^\top X + \lambda I_p)^{-1} X^\top y\)

Properties
If \(\lambda\) is close to zero, the solution \(\hat{\beta}_{\text{PLS}}\) will be close to the classical least squares estimator.
If \(\lambda\) is large, the penalty term dominates strongly and the solution \(\hat{\beta}_{\text{PLS}}\) will be close to 0.

Penalization of \(\beta_0\) is often not desired. This can be avoided by setting
\(\text{pen}(\beta) = \beta^\top K \beta\)
with \(K = \text{diag}(0, 1, 1, \dots, 1)\).
This approach can be generalized be choosing \(K\) as some other positive semidefinite matrix.

The solution is
\(\hat{\beta}_{\text{PLS}} = (X^\top X + \lambda K)^{-1} X^\top y\)
with expectation vector and covariance matrix
\(\mathbb{E}(\hat{\beta}_{\text{PLS}}) = (X^\top X + \lambda K)^{-1} X^\top X \beta\)
\(\text{Cov}(\hat{\beta}_{\text{PLS}}) = \hat{\sigma}^2 (X^\top X + \lambda K)^{-1} X^\top X (X^\top X + \lambda K)^{-1}\)

Ridge regression is associated with bias and variance reduction in comparison to classical least squares. 
An appropriate choice of \(\lambda\) is key when seek a good balance. For this purpose, e.g. cross validation methods may be used.


LASSO = least absolute shrinkage and selection operator

Motivation
Ridge regression hardly penalizes small values, but penalizing and setting small values to zero would reduce complexity.
This would conveniently combine estimation and variable selection.

Leaving out the intercept \(\beta_0\), LASSO uses as penalty the absolute values:
\(\text{pen}(\beta) = \sum_{j = 1}^k |\beta_j|.\)

The LASSO estimator \(\hat{\beta}_{\text{LASSO}}\) minimizes
\((y - X \beta)^\top (y - X \beta) + \lambda \cdot \sum_{j = 1}^k |\beta_j|.\)



Define time series and their applications

Time series play an important role in many application domains:
Price evolutions in financial markets, e.g., prices of stocks, bond, indices,...
- Unemployment rates
- Supply and demand of products
- Insurance losses over time
- Temperatures and other weather data

A time series describes the evolution of variables over time, e.g., oberservations \( x_1, x_2, \dots, x_t \) at discrete
time points \( 1, 2, \dots, t \), which could be modelled as r.v. and/or a discrete stochastic process.



Describe the basic building blocks of time series and stationarity

For a discrete stochastic process \( (X_t)_{t \in \mathbb{Z}} \) in \( (\Omega, \Sigma, \mathbb{P}) \)
we have the marginal distributions

\( F_{t_1, \dots, t_n}(x_1, \dots, x_n) = \mathbb{P}(X_{t_{1}} \leq x_1, \dots, X_{t_{n}} \leq x_n) \)
\( F_{t_n \mid t_{n-1}, \dots, t_1}(x_n \mid x_{n-1}, \dots, x_1) = \mathbb{P}(X_{t_{n}} \leq x_n \mid X_{t_{n-1}} \leq x_{n-1}, \dots X_1 \leq x_{1}) \)

which induces the associated mean and autocovariance functions

\( \mu_t = \mathbb{E}[X_t] = \int xdF_t(x), t \in \mathbb{Z} \)
\( \gamma(t, \tau) = \mathbb{E}[(X_t - \mu_t)(X_{t-\tau} - \mu_{t-\tau}] = \int (x - \mu_t)(y - \mu_{t-\tau})dF_{t,t-\tau}(x,y), t \in \mathbb{Z} \)

For a stochastic process \(X = (X_t)_{t \in \mathbb{Z}}\) two different notions of stationarity are most important in the context of time series:

Covariance stationarity: The mean and covariance function do not depend on time, i.e.,
    \(\forall t, \tau: \mu_t = \mu\)
    \(\forall t, \tau: \gamma(t, t + \tau) = \gamma_\tau\)
Strict stationarity: \(\forall n, \forall t_1, t_2, \dots, t_n, s: (X_{t_1}, X_{t_2}, \dots, X_{t_n}) \stackrel{d}{=} (X_{t_1 + s}, X_{t_2 + s}, \dots, X_{t_n + s})\)

Remark
Covariance stationarity is often called weak stationarity.
However, strict stationarity implies weak stationarity, iff the covariance exists, i.e., if the stochastic process lives in \(L^2\).

The autocorrelation function (ACF) of a weakly stationary process is then \( \rho(\tau) = \frac{\gamma(t)}{\gamma(0)} \in [−1, 1]\).

White noise: A weakly stationary process is called white noise, if the following properties hold:
\(\forall t: \mu_t = 0\)
\(\forall t: \gamma_t = \sigma^2 \cdot \delta_0\)
Independent white noise: White noise \(X\) is called strict white noise, if \(X_1, X_2, \dots\) are iid.
Random walk: A stochastic process \(X\) that can be represented by 
\(X_t = c + X_{t - 1} + \epsilon_t\) with a constant \(c\) and white noise \(\epsilon\) is called a random walk. The constant \(c\) is called the drift.

Example

The random walk \(X_t = \epsilon_t + \epsilon_{t - 1} + \dots + \epsilon_1\) has mean function \(\mu_t = 0\) 
and autocovariance function \(\gamma(t, \tau) = (t - \tau) \sigma^2\) with \(\tau \le t\). 
The autocorrelation function is \(\rho(t, \tau) = \sqrt{1 - \frac{\tau}{t}}\). The random walk is not stationary.


Define an \( AR(1) \) model

The stochastic process \(X\) is an autoregressive process of first order, i.e., an AR(1) process, if
\(X_t = c + \alpha X_{t - 1} + \epsilon_t\)
with a constant parameter \(\alpha, |\alpha| < 1\).
The process is indexed by time \(t \in \mathbb{Z}\). For each \(t\), one assumes that the processes comes from \(-\infty\). This creates a strictly stationary process for strict white noise \((\epsilon_t)_t\).
Exploiting the recursive structure of autoregressive processes, one can write

\(X_t = c \cdot \frac{1 - \alpha^k}{1 - \alpha} + \alpha^k \cdot X_{t - k} + \sum_{i = 0}^{k - 1} \alpha^i \cdot \epsilon_{t - i}\)

If \(X\) is strictly stationary, \(\alpha^k \cdot X_{t - k} \rightarrow 0\) if \(k \rightarrow \infty\). Setting \(\mu_t = c / (1 - \alpha)\), \(\gamma_\tau = \sigma_\epsilon^2 \alpha^2 / (1 - \alpha^2)\) and \(\rho_\tau = \alpha^\tau\) this leads to the solution

\(X_t = c \cdot \frac{1}{1 - \alpha} + \sum_{i = 0}^\infty \alpha^i \cdot \epsilon_{t - i}\)


Give methods to estimate the trend

Linear models of the trend:
Linear trend: \(T_t = \alpha_0 + \alpha_1 t\)
Polynomial trend: \(T_t = \alpha_0 + \alpha_1 t + \alpha_2 t^2 + \dots + \alpha_k t^k\)
Nonlinear models of the trend:
Exponential trend: \(T_t = \alpha \exp(\beta t)\)
Saturation model: \(T_t = \frac{\alpha}{1 + \beta \exp(-\gamma t)}\)

Moving windows:
Let \(w_j \ge 0, j = -q, \dots, q, q \in \mathbb{N}\) be positive weights with \(\sum_j w_j = 1\).
We consider for observations \(x_1, \dots, x_n\) (\(n\) large) the moving window estimator of length \(2q + 1\):

\(\hat{T}_t := \sum_{j = -q}^q w_j x_{t - j}\)


We consider a linear trend
\(T_t = \alpha_0 + \alpha_1 t\)

The corresponding least squares optimization problem with controls \(\alpha_0, \alpha_1\) is given by:
\[\sum_{t = 1}^n (x_t - \alpha_0 - \alpha_1 t)^2 = \min_{\alpha_0, \alpha_1}\]
The explicit solution is
\[\hat{\alpha}_1 = \frac{\sum_{t = 1}^n (x_t - \bar{x})(t - \bar{t})}{\sum_{t = 1}^n (t - \bar{t})^2}, \quad \hat{\alpha}_0 = \bar{x} - \hat{\alpha}_1 \bar{t}\]
with \(\bar{x} = \frac{1}{n} \sum_{t = 1}^n x_t, \bar{t} = \frac{1}{n} \sum_{t = 1}^n t\).

This provides the plug-in estimator of the trend:
\[\hat{T}_t = \hat{\alpha}_0 + \hat{\alpha}_1 t\]


Define linear processes and state Wold's theorem

A linear process admits a representation of the following form:
\(X_t = \mu + \sum_{i = -\infty}^\infty a_i \epsilon_{t - i}\)

with white noise \(\epsilon\) with variance \(\sigma^2\) and \(\sum_{i = -\infty}^\infty |a_i| < \infty\).

A linear process is weakly in \(L^2\) stationary with \(\mu_t = \mu\) and \(\gamma(t, \tau) = \sigma^2 \sum_{i = -\infty}^\infty a_i a_{i - \tau}\)
In fact, by Wold's Decomposition Theorem, any weakly stationary time series that is purely non deterministic can be represented as a linear process of the form

\(X_t = \mu + \sum_{i = 0}^\infty a_i \epsilon_{t - i}\)
\(a_0 = 1, \sum_{i = 1}^\infty a_i^2 < \infty\)

While stationary processes are key building blocks, many processes are non-stationary. Like a random walk they may be sums of stationary processes. Generalizing this idea leads to integrated processes.


Theorem (Wold's Decomposition Theorem): 
Any zero-mean nondeterministic covariance-stationary process \(\{x_t; t \in \mathbb{Z}\}\) can be decomposed as
\(x_t = \sum_{j = 0}^\infty b_j \varepsilon{t - j} + d_t\)
where
\(b_0 = 1, \sum_{j = 1}^\infty b_j^2 < \infty\),
\(\varepsilon_t \sim WN(0, \sigma_u^2)\), (white noise)
\(\{b_j\}\) and \(\{\varepsilon_t\}\) are unique,
\(\{d_t; t \in \mathbb{Z}\}\) is deterministic,
\(\varepsilon_t\) is the limit of linear combinations of \(X_s, s \le t\)
\(E(d_t \varepsilons_s) = 0 \forall t, s\)


Define integrated processes


The lag operator \(L\) shifts the process in time, i.e.,
\(LX_t := (LX)_t := X_{t - 1}\)

The difference operator is defined as
\(\Delta := \text{id} - L, \Delta^k = (\text{id} - L)^k\)

Integrated process
A process \(X\) is integrated of order \(d\), \(I(d)\), if

\(\Delta^{d - 1} X\) is non-stationary, and
\(\Delta^d X\) is stationary.

White noise and an AR(1) process are examples of \(I(0)\) processes, a random walk is \(I(1)\). In practice, \(I(d)\) processes for \(d > 1\) are less frequently used.
In the following, we observe that the observed process \(Y\) is \(I(d)\), but we analyze the transformed stationary process \(X = \Delta^d Y\).

Define general \( \text{MA}(q) \) processes and prove their stationarity

A moving average process of order \(q\), \(MA(q)\), is a process of the following form
\(X_t = \beta_1 \epsilon_{t - 1} + ... + \beta_q \epsilon_{t - q} + \epsilon_t\)

with white noise \(\epsilon\).
An \(MA(q)\) process can conveniently be rewritten using the Lag operator:

\(X_t = \beta(L) \epsilon_t\)
\(\beta(L) = 1 + \beta_1 L + ... + \beta_q L^q\)

Interpretation
White noise models shocks at different points in time \(t\).
The \(MA(q)\) process is a weighted average of shocks at different times. 
When considering the processes at a different point in time, the averaging window is moved to the new position.

\(X_t = \beta(L) \epsilon_t,\)
\(\beta(L) = 1 + \beta_1 L + ... + \beta_q L^q, \beta_0 = 1, \text{Var}(\epsilon_t) = \sigma^2\)

is stationary with
\(\mu_t = 0, \gamma_\tau := \gamma(t, \tau) = \begin{cases} \sum_{i = 0}^{q - |\tau|} \beta_i \beta_{i + |\tau|} \cdot \sigma^2; & |\tau| \le q \\ 0 & |\tau| > q \end{cases}\)
\(\rho_\tau = \begin{cases} \frac{\sum_{i = 0}^{q - |\tau|} \beta_i \beta_{i + |\tau|}}{\sum_{i = 0}^q \beta_i^2} & |\tau| \le q \\ 0 & |\tau| > q \end{cases}\)

Example
For the MA(1) process
\(X_t = \beta \epsilon_{t - 1} + \epsilon_t\)
one obtains that \(\rho_\tau = 0\) for \(|\tau| > 1\) and \(\rho_1 = \frac{\beta}{1 + \beta^2}\).


Define the general \( \text{AR}(p) \) process and show its stationarity
The linear autoregressive process of order \(p\), \( \text{AR}(p) \), is defined as

\(X_t = \nu + \alpha_1 X_{t - 1} + ... + \alpha_p X_{t - p} + \epsilon_t\)
with white noise \(\epsilon\).

An AR(\(p\)) process can conveniently be rewritten using the Lag operator:
\(\alpha(L)X_t = \nu + \epsilon_t\)
\(\alpha(L) = 1 - \alpha_1 L + ... - \alpha_p L^p\)

Interpretation
The defining equation is formally a multiple regression.
The explanatory variables are the historic realisations of the process.


The AR(p) process

\(\alpha(L)X_t = \nu + \epsilon_t\)
\(\alpha(L) = 1 - \alpha_1 L + ... - \alpha_p L^p\)

is stationary iff.

\(\forall z \in \mathbb{C}, |z| \le 1: \alpha(z) \ne 0\) with \(\alpha(z) = 1 - \alpha_1 z + ... - \alpha_p z^p\),
i.e., all complex roots of \(\alpha\) lie strictly outside the unit circle.

Representation as MA(\(\infty\)) process
In the stationary case of the AR(\(p\)) process one determines the series \(\frac{1}{\alpha(z)} = \sum_{i = 0}^\infty a_i z^i\) which leads to the representation

\(X_t = \sum_{i = 0}^\infty a_i \nu + \sum_{i = 0}^\infty a_i L^i \epsilon_t\)


State and derive the Yule-Walker equations

Consider a stationary AR(\(p\)) process with \(\nu = 0: \alpha(L)X_t = \epsilon_t, \alpha(L) = 1 - \alpha_1 L + ... - \alpha_p L^p\). 
The Yule-Walker equations describe a simple correspondence between \(\alpha = (\alpha_1, \alpha_2, ..., \alpha_p)\) 
and the autocorrelations \(\rho = (\rho_1, \rho_2, ..., \rho_p)\).

These can be derived by observing for \(\tau = 1, 2, ..., p\) 
that \(E(X_t X_{t - \tau}) = \alpha_1 E(X_{t - 1} X_{t - \tau}) + \dots + \alpha_p E(X_{t - p} X_{t - \tau}), \tau = \gamma_{-\tau}\)
Define the matrix, which is Toepelitz

\(R = \begin{pmatrix} 1 & \rho_1 & \rho_2 & \rho_3 & ... & \rho_{p - 1} \\ \rho_1 & 1 & \rho_1 & \rho_2 & ... & \rho_{p - 2} \\ \rho_2 & \rho_1 & 1 & \rho_1 & ... & \rho_{p - 3} \\ ... & ... & ... & ... & ... & ... \\ \rho_{p - 1} & \rho_{p - 2} & \rho_{p - 3} & \rho_{p - 4} & ... & 1 \end{pmatrix}\)

Yule-Walker equations
The coefficients \(\alpha\) and autocorrelations \(\rho\) can be transformed into each other by the Yule-Walker equations:

\(\rho = R \alpha\)
Moreover, if \(\tau > p\), then \(\rho_\tau = \alpha_1 \rho_{\tau - 1} + \alpha_2 \rho_{\tau - 2} + ... + \alpha_p \rho_{\tau - p}\).


Define ARMA processes and state the requirements for stationarity


Centered ARMA
A stochastic process \(X\) is a centered ARMA(\(p, q\)) process, if it satisfied the following equation:

\(X_t - \alpha_1 X_{t - 1} - \dots - \alpha_p X_{t - p} = \epsilon_t + \beta_1 \epsilon_{t - 1} + \dots + \beta_q \epsilon_{t - q}\)

The terms \(\alpha_1, \dots, \alpha_p, \beta_1, \dots, \beta_q\) are real parameters and \((\epsilon_t)\) is white noise.
ARMA with mean \(\mu\)

The process \(X\) is an ARMA(\(p, q\)) process with mean \(\mu\) if \((X_t - \mu)\) is a centered ARMA(\(p, q\)) process.
Characteristic polynomials

We define the characteristic polynomials of the process by
\(\beta(z) = 1 + \beta_1 z + \dots + \beta_q z^q\)
\(\alpha(z) = 1 - \alpha_1 z - \dots - \alpha_p z^p\)

The centered ARMA(\(p, q\)) can then be represented by

\(\alpha(L)X_t = \beta(L) \epsilon_t\)

Common roots of \(\alpha\) and \(\beta\) are excluded for technical reasons.
ARMA(\(p, q\)): \(\alpha(L)X_t = \beta(L) \epsilon_t\)
Stationarity:
The process is stationary, iff the roots of \(z \rightarrow \alpha(z)\) lie strictly outside the unit circle.
In this case, the centered ARMA process possesses an MA(\(\infty\))-representation:

\(X_t = \alpha(L)^{-1} \beta(L) \epsilon_t\)

Invertibility:
The process is called invertible, iff the roots of \(z \rightarrow \beta(z)\) lie strictly outside the unit circle.
In this case, the centered ARMA process possesses an AR(\(\infty\))-representation:
\(\beta^{-1}(L) \alpha(L) X_t = \epsilon_t\)

Approximation:
Every stationary, invertible ARMA(\(p, q\)) process can be represented as both an MA(\(\infty\)) process and an AR(\(\infty\)).
This implies that it can be approximated by both a pure MA and AR process of sufficiently high order.
Methods for AR and MA are thus always applicable.
In addition, Wold's theorem shows that ARMA processes with arbitrary mean, 
provide approximate representations of all weakly stationary processes that are purely non deterministic.


Define ARIMA processes


Difference operator:
\(\Delta Y_t = Y_t - Y_{t - 1}\) with \(\Delta^d Y_t = \begin{cases} \Delta Y_t, & d = 1, \\ \Delta^{d - 1} (\Delta Y_t), & d > 1. \end{cases}\)

ARMA(\(p, q\))
A time series \(Y\) is called ARIMA(\(p, q, d\)) process (autoregressive integrated moving average), if the differenced process \(X = \Delta^d (Y)\) is an ARMA(\(p, q\)) process.
Non stationarity
For \(d \ge 1\) ARIMA processes are non stationary.

Application to Time Series Data

Differencing (once or multiple time) may transform a non-stationary data set into one that may be described or modeled by an ARMA(\(p, q\)) process.
An example are log-returns modeled by an ARMA(\(p, q\)) process that lead to ARIMA(\(p, q, 1\)) time series for the corresponding log prices (\(S_t\)).



State the estimators for the statistical properties of time series

Estimator of the mean µ = E (Xt ) is the arithmetic sample mean
\( \bar{X}_n = \frac{1}{n}\sum_{i=1}^n X_i \)
its unbiased and has variance, where \( \tau \) is the covariance function
\( \text{Var}(\bar{X}_n) = \frac{1}{n}\sum_{\tau=-(n-1)}^{n-1} \frac{n - \tau}{n}\gamma_\tau < \infty \)
which is in the limit
\( \lim_{n \to \infty} n \text{Var}(\bar{X}_n) = \gamma_0 + 2\sum_{\tau = 1}^\infty \gamma_\tau \coloneqq f(0) \)
and we have the asymptotic result
\( \sqrt{n}(\bar{X}_n - \mu) \xrightarrow{d} N(0, f(0)) \)

The estimator for the covariance function is
\( \hat{\gamma}_{\tau,n} = \frac{1}{n}\sum_{t=1}^{n-\tau} (X_t - \bar{X}_n)(X_{t + \tau} - \bar{X}_n)\)
this estimator is not unbiased, nevertheless asymptotically it holds 
\( \sqrt{n}(\bar{\gamma}_{\tau,n} - \gamma_\tau) \xrightarrow{d} N(0, \sigma^2_{\tau, \infty}) \)
where \( \sigma^2_{\tau, \infty} = \sum^{\infty}_{-\infty} (\gamma^2_j + \gamma_{j - \tau}\gamma_{j + \tau}\).

For the autocorrelation function we have the estimator

\( \hat{\rho}_{\tau,n} = \frac{\hat{\gamma}_{\tau,n}}{\hat{\gamma}_{0,n}} \)
which again has bias, namely \( \mathbb{E}[\hat{\rho}_{\tau,n}] = \rho_\tau + \mathcal{O}(n^{-1}) \)
the corresponding variance contains a rather garstly term which we shorten to \( \Sigma_{\rho,\tau \tau} \) 
and get \( \text{Var}(\hat{\rho}_{\tau, n} = \frac{1}{n} \Sigma_{\rho,\tau \tau} + \mathcal{O}(n^{-2})\)
and asymptotically it holds 
\( \sqrt{n}(\rho_{(m),n} - \rho_{(m)}) \xrightarrow{d} N(0, \Sigma_\rho) \)


We recall that \(\sqrt{n} (\hat{\rho}_{\tau, n} - \rho_\tau) \stackrel{d}{\longrightarrow} N(0, \Sigma_{\rho, \tau \tau})\).
Since \(\rho_\tau = 0\) for all \(\tau > q\), the formula for the asymptotic variance may be simplified

\(\Sigma_{\rho, \tau \tau} = 1 + 2 \sum_{i = 1}^q \rho_i^2\)

White noise:
If \(X\) is white noise, then \(E[\hat{\rho}_{\tau, n}] = -n^{-1} + (n - 2)\) and \(\Sigma_\rho\) is the identity matrix.
This asymptotic results may be exploited to construct asymptotic confidence intervals.


Describe the estimation of \( \text{AR}(p) \) processes

The AR(\(p\))-model satisfies the following equation
\(X_t = \alpha_1 X_{t - 1} + ... + \alpha_p X_{t - p} + \epsilon_t\)
with Var(\(\epsilon_t\)) \(= \sigma^2\).

We already know and analyzed a method to estimate the ACF.
The Yule-Walker estimator is a plug-in estimator that computes the estimated parameters \(\hat{\alpha}_1, \hat{\alpha}_2, ..., \hat{\alpha}_p\) from the Yule-Walker equations:

\(\begin{pmatrix} 1 & \hat{\rho}_1 & \hat{\rho}_2 & \hat{\rho}_3 & ... & \hat{\rho}_{p - 1} \\ \hat{\rho}_1 & 1 & \hat{\rho}_1 & \hat{\rho}_2 & ... & \hat{\rho}_{p - 2} \\ \hat{\rho}_2 & \hat{\rho}_1 & 1 & \hat{\rho}_1 & ... & \hat{\rho}_{p - 3} \\ ... & ... & ... & ... & ... & ... \\ \hat{\rho}_{p - 1} & \hat{\rho}_{p - 2} & \hat{\rho}_{p - 3} & \hat{\rho}_{p - 4} & ... & 1 \end{pmatrix} \begin{pmatrix} \hat{\alpha}_1 \\ \hat{\alpha}_2 \\ \hat{\alpha}_3 \\ ... \\ \hat{\alpha}_p \end{pmatrix} = \begin{pmatrix} \hat{\rho}_1 \\ \hat{\rho}_2 \\ \hat{\rho}_3 \\ ... \\ \hat{\rho}_p \end{pmatrix}\)

Alternatively, the model may also be interpreted as a linear model, and the estimated parameters
\( \hat{\alpha}_1, \hat{\alpha}_2, \dots, \hat{\alpha}_p \) may be obtained from least squares estimation.

The Yule-Walker estimator and the Least Squares estimator are asymptotically equivalent for Gaussian error terms.
In this case, the estimators are consistent and asymptotically normal with covariance matrix

\(\sigma^2 \Gamma^{-1}\) with \(\Gamma = \begin{pmatrix} 1 & \gamma_1 & \gamma_2 & \gamma_3 & ... & \gamma_{p - 1} \\ \gamma_1 & 1 & \gamma_1 & \gamma_2 & ... & \gamma_{p - 2} \\ \gamma_2 & \gamma_1 & 1 & \gamma_1 & ... & \gamma_{p - 3} \\ ... & ... & ... & ... & ... & ... \\ \gamma_{p - 1} & \gamma_{p - 2} & \gamma_{p - 3} & \gamma_{p - 4} & ... & 1 \end{pmatrix}\)


Define ARCH processes

Suppose that \((Z_t)_{t \in \mathbb{Z}}\) is strict white noise with variance 1, i.e., SWN(0, 1). 
The process \((\epsilon_t)_{t \in \mathbb{Z}}\) is called an ARCH(\(p\)) process (autoregressive conditionally heteroscedastic), if:

\((\epsilon_t)_{t \in \mathbb{Z}}\) is strictly stationary.
For some strictly positive process \((\sigma_t)_{t \in \mathbb{Z}}\) the following identities hold:

\(\epsilon_t = \sigma_t Z_t\)
\(\sigma_t^2 = \alpha_0 + \sum_{i = 1}^p \alpha_i \epsilon_{t - i}^2\)

with constants \(\alpha_0 > 0, \alpha_1, ..., \alpha_p \ge 0\).
Stationarity and martingale difference sequence

An ARCH(\(p\)) process is a martingale difference sequence with respect to its natural filtration:
\(E(\epsilon_t | F_{t - 1}) = E(\sigma_t Z_t | F_{t - 1}) = \sigma_t E(Z_t | F_{t - 1}) = \sigma_t E(Z_t) = 0\)

If an ARCH(\(p\)) process is in \(L^2\), it is weakly stationary. This requires \(\alpha_1 + ... + \alpha_p < 1\).
A covariance stationary ARCH(\(p\)) process is white noise.


Stochastic Recursion
\(\epsilon_t = \sqrt{\alpha_0 + \sum_{i = 1}^p \alpha_i \epsilon_{t - i}^2} \cdot Z_t\)
Typical examples for \((Z_t)_{t \in \mathbb{Z}}\):

a) \(Z_t \sim N(0, 1)\) (normal distribution)
b) \(Z_t \sim t_\nu\) (\(t\)-distribution)

Distribution:
The distribution of \(\epsilon_t\) is typically not characterized by simple standard distributions.
In case a), one obtains a normal variance mixture.


Conditional expectation and volatility:

Letting \(\mathcal{F}_t = \sigma(\epsilon_s: s \le t)\) be the \(\sigma\)-algebra of the past up to time \(t\), 
\(t \in \mathbb{Z}\), we obtain the natural filtration of the process \(\epsilon\).

Conditional expection: \(E[\epsilon_t | \mathcal{F}_{t - 1}] = 0\), i.e., \(\epsilon\) is a martingale difference sequence. 
If \(\epsilon\) models return, the expected return conditional on the available information is 0.
Volatility: \(\sqrt{\text{Var}(\epsilon_t | \mathcal{F}_{t - 1})} = \sigma_t\)
- Heteroscedasticity: Volatility/conditional variance vary randomly over time.
- Volatility cluster: If \(\epsilon_{t - 1}^2, \epsilon_{t - 2}^2, ...\) are large, then \(\sigma_t^2\) is large.
